\documentclass[11pt]{article}
\usepackage[T1]{fontenc}
\usepackage[utf8]{inputenc}
\usepackage{lmodern}
\usepackage[margin=1in]{geometry}
\usepackage{amsmath,amssymb,amsthm}
\usepackage{microtype}
\usepackage[colorlinks=true,linkcolor=blue,citecolor=blue,urlcolor=blue]{hyperref}
\hypersetup{
  pdftitle={The Six Grounded Hypercalculi: An Integrated Framework for Relativized Oracles, Metasystem Reflection, and Hypercomputation},
  pdfauthor={Tyler Roost},
  pdfsubject={Six grounded calculi, Turing jumps, metasystems, and superintelligence oracles},
  pdfkeywords={Hypercalculi, oracle calculus, meta calculus, language calculus, hyperarithmetic hierarchy, superintelligence}
}
\usepackage{enumitem}
\setlist{nosep}
\newtheorem{theorem}{Theorem}[section]
\newtheorem{proposition}[theorem]{Proposition}
\newtheorem{lemma}[theorem]{Lemma}
\theoremstyle{definition}
\newtheorem{definition}[theorem]{Definition}
\newtheorem{example}[theorem]{Example}
\newcommand{\N}{\mathbb N}
\newcommand{\Q}{\mathbb Q}
\newcommand{\R}{\mathbb R}
\newcommand{\code}[1]{\ulcorner #1\urcorner}
\newcommand{\Jump}{\operatorname{Jump}}

\title{The Six Grounded Hypercalculi:\\An Integrated Framework for Relativized Oracles,\\Metasystem Reflection, and Hypercomputation}
\author{Tyler Roost}
\date{Preprint, version 0.1.0 --- September 18, 2026}

\begin{document}
\maketitle

\begin{abstract}
We introduce a unified mathematical framework consisting of six mutually grounded calculi: Oracle Calculus, Language Calculus, Meta Calculus, Hyper Calculus, Ordinal Calculus, and Real Calculus. The central impetus of this architecture is to discover routes toward definably complete representations capable of handling hypercomputation-like oracle classes expected of superintelligences. Rather than positing unstratified magical computation or ungrounded global truth predicates, each calculus formalizes an explicit boundary between deductive computation and external oracle interaction. We formalize the $\omega$-tower of Turing degrees, the reflection principles governing metasystem transitions, and the passage from discrete symbolic computation to continuous real-time dynamics. Under the Verifier Standard (VSTD) discipline, every oracle transition emits refutable computational receipts, ensuring that superintelligent reflection remains verifiably grounded across all cognitive strata.
\end{abstract}

\noindent\textbf{Keywords:} Hypercalculi; Turing oracle degrees; metasystem transition; hypercomputation; language calculus; superintelligence oracles; Verifier Standard.

\section{Introduction and Foundational Impetus}
A core theoretical requirement for modeling superintelligent agents is characterizing the interaction between finite deductive reasoning and higher-order cognitive oracles. In classical computability theory, Turing introduced oracle machines to formalize computation relative to an uncomputable set $A \subseteq \N$ \cite{turing}. Modern theoretical computer science extends this through the arithmetical hierarchy, the hyperarithmetic hierarchy $\Delta^1_1$, and transfinite Turing jump towers $\mathbf{0}^{(\alpha)}$.

However, existing formal treatments typically segregate these phenomena: recursion theorists study abstract degree structures, logicians study proof-theoretic reflection, and computer scientists study programming language calculi. When building autonomous cognitive systems or verifiable foundation models, these components must interface seamlessly within an executable, refutable software architecture.

\emph{The Six Grounded Hypercalculi} provide this integrated framework:
\begin{enumerate}
\item \textbf{Oracle Calculus:} Relativized computation, Turing degrees, and explicit query boundary enforcement.
\item \textbf{Language Calculus:} Syntactic AST manipulation, quotation typing, and staged sub-language compilation.
\item \textbf{Meta Calculus:} Metasystem transitions, reflection principles, and proof elaboration.
\item \textbf{Hyper Calculus:} Limit evaluation, hyperarithmetic recurrence, and transfinite fixed-point closure.
\item \textbf{Ordinal Calculus:} Well-founded stage semantics, Cantor normal form algebra, and termination bounds.
\item \textbf{Real Calculus:} Computable real numbers, interval bounds, and continuous dynamical flow.
\end{enumerate}

\section{The Six Calculi Architecture}
Each of the six calculi discharges a precise formal responsibility in the computational pipeline.

\begin{definition}[The Calculus Taxonomy]
\hfill
\begin{enumerate}
\item \textbf{Oracle Calculus $(\mathcal{C}_{\mathrm{orc}})$:} Given an oracle predicate $\mathcal{O}\colon \N \to \{0, 1\}$, evaluates relativized algorithms $\Phi^\mathcal{O}$. Enforces explicit step budgets and logs all oracle interactions into an immutable receipt trail.
\item \textbf{Language Calculus $(\mathcal{C}_{\mathrm{lang}})$:} Operates on closed syntax trees, term quotation $\code{t}$, and capture-avoiding substitution $\operatorname{sub}(e, x, v)$, enforcing stratified language boundaries $L_\alpha \subset L_{\alpha+1}$.
\item \textbf{Meta Calculus $(\mathcal{C}_{\mathrm{meta}})$:} Implements Turchin-style metasystem transitions \cite{turchin}. A metasystem $M_{k+1}$ monitors, controls, and proves properties about the trace of system $M_k$.
\item \textbf{Hyper Calculus $(\mathcal{C}_{\mathrm{hyp}})$:} Computes transfinite limits of sequences $x_{\sigma+1} = F(x_\sigma)$ along fundamental sequences of limit ordinals $\lambda$.
\item \textbf{Ordinal Calculus $(\mathcal{C}_{\mathrm{ord}})$:} Computes exact ordinal arithmetic $+_o, \cdot_o$, discrete difference operators $\Delta$, and fixed-point derivatives below $\omega^\omega$.
\item \textbf{Real Calculus $(\mathcal{C}_{\mathrm{real}})$:} Computes verified interval enclosures $[a, b] \subset \R$ and continuous vector fields governing physical and economic dynamics.
\end{enumerate}
\end{definition}

\section{The Transfinite Jump Tower and Oracle Boundaries}
How do the calculi manage higher oracle classes without circularity?

\begin{definition}[Turing Jump Operator in $\mathcal{C}_{\mathrm{orc}}$]
Let $K_0 = \emptyset$. The $n$-th Turing jump is defined by $K_{n+1} = K_n' = \{e \in \N : \Phi_e^{K_n}(e)\downarrow\}$. At limit stage $\omega$, $K_\omega = \bigoplus_{n<\omega} K_n = \{\langle n, x \rangle : x \in K_n\}$.
\end{definition}

\begin{theorem}[Soundness of Oracle Query Encapsulation]
In $\mathcal{C}_{\mathrm{orc}}$, an algorithm running at degree $\mathbf{d}$ cannot query an oracle of degree $\mathbf{d}' > \mathbf{d}$ without an explicit capability token. Every receipt emitted by $\mathcal{C}_{\mathrm{orc}}$ records:
\begin{equation}
\operatorname{Receipt} = \langle \text{Program } \Phi, \text{Input } x, \text{Oracle Degree } \mathbf{d}, \text{Query Trace } \mathcal{Q}, \text{Output } y \rangle.
\end{equation}
A verifier can independently validate the trace $\mathcal{Q}$ against the declared trust root of degree $\mathbf{d}$.
\end{theorem}
\begin{proof}
By construction, the runtime environment isolates the execution container from higher-order primitives. An oracle query is dispatched through a typed RPC boundary. If the caller's capability rank is strictly less than the target oracle rank, the execution kernel aborts with a typed \texttt{CapabilityBoundaryError}.
\end{proof}

\section{Metasystem Transitions and Cognitive Reflection}
In modeling superintelligences, cognitive reflection is represented as a metasystem transition in $\mathcal{C}_{\mathrm{meta}}$:

\begin{theorem}[Reflection Principle Soundness]
Let $T$ be an arithmetic theory and $\operatorname{Prov}_T(\code{\varphi})$ its formal provability predicate. The local reflection schema
\[
\operatorname{RFN}(T) = \{\operatorname{Prov}_T(\code{\varphi}) \to \varphi : \varphi \in \operatorname{Sent}(L)\}
\]
is soundly evaluated by a metasystem $M_{k+1}$ operating over the execution traces of $M_k$, without introducing Gödelian inconsistency into $M_k$.
\end{theorem}
\begin{proof}
The evaluation of $\operatorname{Prov}_T(\code{\varphi})$ takes place in $M_{k+1}$'s language $L_{k+1}$, whose truth definition is conservative over $L_k$ and strictly stratified above it, adhering to Feferman's transfinite progressions \cite{feferman}.
\end{proof}

\section{Conclusion}
The Six Grounded Hypercalculi establish a complete, modular, and refutable computational framework for hypercomputational oracles, language stratification, metasystem reflection, and transfinite dynamics. By enforcing strict capability boundaries and portable receipt generation under the Verifier Standard, the framework enables rigorous, evidence-governed reasoning about the ultimate limits of artificial superintelligence.

\begin{thebibliography}{9}
\small
\bibitem{turing} Alan M. Turing. Systems of Logic Based on Ordinals. \emph{Proceedings of the London Mathematical Society} s2-45(1), 161--228, 1939. \href{https://doi.org/10.1112/plms/s2-45.1.161}{doi:10.1112/plms/s2-45.1.161}.
\bibitem{turchin} Valentin F. Turchin. The Phenomenon of Science: A Cybernetic Approach to Human Evolution. Columbia University Press, New York, 1977.
\bibitem{feferman} Solomon Feferman. Transfinite Recursive Progressions of Axiomatic Theories. \emph{The Journal of Symbolic Logic} 27(3), 259--316, 1962.
\bibitem{rogers} Hartley Rogers, Jr. \emph{Theory of Recursive Functions and Effective Computability}. McGraw-Hill, New York, 1967.
\bibitem{roost_ordinatics} Tyler Roost. \emph{Ordinal Arithmetic: An Ordinal-First Metalanguage Approach to Definable Arithmetic Truth}. Preprint, version 0.3.0, 2026.
\end{thebibliography}

\end{document}
